\documentclass[12pt,a4paper]{report}

% Limbă și diacritice
\usepackage{fontspec}
\defaultfontfeatures{Ligatures=TeX}
\setmainfont{Noto Serif}

% Margini & font
\usepackage[margin=2.5cm]{geometry}
\usepackage{lmodern}
\linespread{1.1}

% Hiperlinkuri
\usepackage[hidelinks]{hyperref}

% Tabele & figuri
\usepackage{graphicx}
\usepackage{longtable}
\usepackage{float}
\usepackage{booktabs}
% \usepackage{siunitx}
% \sisetup{output-decimal-marker={,}}
\usepackage{subcaption}

% Anexe
\usepackage[toc,page]{appendix}
\renewcommand{\appendixpagename}{Anexe}
\renewcommand{\appendixtocname}{Anexe}

% Matematică
\usepackage{amsmath, amssymb, amsthm}

% Enumerări
\usepackage{enumitem}

% (Opțiunea 2 - minted, comentată implicit; necesită -shell-escape)
\usepackage{minted}
%\setminted{fontsize=\small, breaklines=true, frame=none}
\usemintedstyle{borland}

% Numerotări: Section -> I., Subsection -> 1.1 etc.
\renewcommand\thesection{\Roman{section}}
\setcounter{secnumdepth}{2}
\setcounter{tocdepth}{2}

% Informații titlu (modifică după nevoie)
\newcommand{\Universitate}{UNIVERSITATEA DE STAT DIN MOLDOVA}
\newcommand{\Facultate}{FACULTATEA DE MATEMATICA SI INFORMATICA}
\newcommand{\Departament}{DEPARTAMENTUL INFORMATICĂ}
\newcommand{\TitluProiect}{Preferințe de cafea și consum zilnic}
\newcommand{\Disciplina}{Statistica pentru Știința Datelor}
\newcommand{\Autor}{Ivan Majeru}
\newcommand{\Titular}{Conf. univ. dr. Ludmila Novac}
\newcommand{\An}{2025}

\begin{document}

% Pagina de titlu
\begin{titlepage}
  \centering
  {\Large \Universitate\par}
  {\large \Facultate\par}
  {\large \Departament\par}
  \vspace{2cm}
  {\LARGE Cercetare statistică\par}
  \vspace{0.5cm}
  {\Large \TitluProiect\par}
  \vspace{2.5cm}
  \begin{flushleft}
    \textbf{Autor:} \Autor\par
    \textbf{Titular disciplină:} \Titular\par
    \textbf{Disciplina:} \Disciplina\par
  \end{flushleft}
  \vfill
  {\An\par}
\end{titlepage}

\pagenumbering{roman}
\tableofcontents
\clearpage
\pagenumbering{arabic}

% I. OBSERVAȚIA STATISTICĂ
\section{Observația statistică}

\subsection{Scopul observației statistice}
Scopul cercetării este de a identifica preferințele de cafea (tipuri preferate, loc de achiziție) și de a analiza consumul zilnic (număr de căni) în funcție de caracteristici socio-demografice.

\subsection{Populația statistică}
Populația vizată include persoane adulte (18+) din Republica Moldova, rezidente în mediul urban și rural.

\subsection{Metoda de observație}
Anchetă statistică prin chestionar online/offline. Eșantionare neprobabilistică (de confort) sau probabilistică (dacă este cazul).

\subsection{Eșantionul observat}
Eșantion de 100 de persoane (date sintetice). Vârsta medie: 46{,}2 ani; consum mediu zilnic: 2{,}75 căni; buget mediu lunar: 1209{,}60 MDL.

\subsection{Caracteristici statistice}
Exemplu:
\begin{itemize}[leftmargin=1.2cm]
  \item Vârstă (ani)
  \item Sex (Masculin/Feminin)
  \item Tip localitate (Urban/Rural)
  \item Tip cafea preferat (Espresso, Latte, Cappuccino, Americano, etc.)
  \item Număr căni/zi
  \item Ora consum dominantă (Dimineața/Prânz/Seara)
  \item Buget lunar pentru cafea (MDL)
  \item Loc achiziție (Acasă-Boabe, Cafenea, Oficiu, etc.)
\end{itemize}

\subsection{Scale de măsurare}
\begin{itemize}[leftmargin=1.2cm]
  \item Vârstă, Număr căni/zi, Buget lunar: scară raport.
  \item Sex, Tip localitate, Tip cafea, Loc achiziție: nominal.
  \item Ora consum dominantă: ordinal (dacă o ordonezi), altfel nominal.
\end{itemize}

\subsection{Sursa datelor}
Date primare sintetice generate pentru scop didactic (Jupyter/Python).

\begin{minted}{python}
import pandas as pd
import numpy as np
from datetime import datetime, timedelta

# Set seed pentru reproducibilitate
np.random.seed(42)
# Numar de inregistrări
n = 100

# Generare date sintetice
data = {"ID": range(1, n+1)}

# Caracteristici demografice
data["Varsta"] = np.random.randint(18, 70, n)
data["Sex"] = np.random.choice(["Masculin", "Feminin"], n, p=[0.48, 0.52])
data["Tip_Localitate"] = np.random.choice(["Urban", "Rural"], n, p=[0.65, 0.35])

# Variabile principale
data["Tip_Cafea_Preferat"] = np.random.choice(
    ["Espresso", "Latte", "Americano", "Cappucino", "Decaf", "Alte"], 
    n,  p=[0.25, 0.20, 0.28, 0.20, 0.05, 0.02]
)
data["Numar_Cani_Pe_Zi" = np.random.choice(
    [1, 2, 3, 4, 5],  n,  p=[0.20, 0.35, 0.25, 0.15, 0.05]
)
data["Ora_Consum_Dominanta"] = np.random.choice(
    ["Dimineata", "Pranz", "Seara"], 
    n,  p=[0.65, 0.25, 0.10]
)
data["Buget_Luna_MDL"] = np.random.randint(200, 2000, n)
data["Loc_Achizitie"] = np.random.choice(
    ["Cafenea", "Supermarket", "Oficiu", "Distributie_Automata",
        "Acasa", "Alte"],
    n, p=[0.25, 0.20, 0.20, 0.20, 0.10, 0.05]
)

# Creare DataFrame
df = pd.DataFrame(data)

\end{minted}

\newpage

E nevoie de unele ajustari logice pentru corerenta.

\begin{minted}{python}
# Persoanele mai tinere tind să consume mai mult latte
df.loc[
    (df["Varsta"] < 30) & (np.random.rand(n) < 0.4),
    "Tip_Cafea_Preferat"
] = "Latte"

# Persoanele mai în varsta tind sa consume mai mult decaf
df.loc[
    (df["Varsta"] > 55) & (np.random.rand(n) < 0.3),
    "Tip_Cafea_Preferat"
] = "Decaf"

# Urban tinde sa aiba buget mai mare
urban_mask = df["Tip_Localitate"] == "Urban"
df.loc[urban_mask, "Buget_Luna_MDL"] = (
    df.loc[urban_mask, "Buget_Luna_MDL"] * np.random.uniform(
    1.1, 1.3, sum(urban_mask))).astype(int)

\end{minted}

% II. SISTEMATIZAREA ȘI PREZENTAREA DATELOR
\section{Sistematizarea și prezentarea datelor observate}

\subsection{Tabel descriptiv}
Aici avem tabelul descriptiv cu 10 elemente din setul de date. Datele complete le puteti vedea in Anexa I.

\begin{table}[H]
\centering
\caption{Primele 10 observații din setul de date}
\label{tab:primele5}
\begin{longtable}{rlllrlrl}
\toprule
Vârsta & Sex & Localitate & Tip Cafea & Cani/zi & Perioada & Buget (MDL) & Loc \\
\midrule
\endfirsthead
\toprule
Vârsta & Sex & Localitate & Tip Cafea & Cani/zi & Perioada & Buget (MDL) & Loc \\
\midrule
\endhead
\midrule
\multicolumn{8}{r}{Continued on next page} \\
\midrule
\endfoot
\bottomrule
\endlastfoot
18 & F & Urban & Latte & 5 & Dimineata & 796 & Supermarket \\
42 & F & Rural & Americano & 2 & Dimineata & 497 & Oficiu \\
60 & M & Rural & Decaf & 4 & Seara & 1151 & Oficiu \\
57 & F & Urban & Cappucino & 3 & Dimineata & 276 & Supermarket \\
59 & M & Rural & Latte & 3 & Pranz & 1307 & Supermarket \\
42 & M & Urban & Cappucino & 3 & Seara & 2022 & Supermarket \\
56 & F & Urban & Decaf & 2 & Dimineata & 2407 & Oficiu \\
52 & M & Urban & Espresso & 2 & Dimineata & 2041 & Supermarket \\
20 & M & Urban & Americano & 2 & Dimineata & 1768 & Distributie \\
67 & M & Urban & Espresso & 3 & Pranz & 2025 & Distributie \\
\end{longtable}

\end{table}

\subsection{Clasificarea unităților observate și ponderea pe grupe}
Exemple de tabele frecvențe/ponderi:
\begin{itemize}[leftmargin=1.2cm]
  \item După frecvența consumului (cani/zi grupate)
  \item După tipul de cafea preferat
  \item După locul de achiziție
  \item După sex / localitate
\end{itemize}

\subsubsection{Dupa frecventa consumului (cani/zi grupate)}

\begin{table}[H]
\centering
\caption{Gruparea dupa frecventa consumului (cani/zi grupate)}
\begin{tabular}{lrr}
\toprule
Nr Cani/Zi & Frecventa & Procent \% \\
\midrule
1 & 18 & 18.00 \\
2 & 41 & 41.00 \\
3 & 30 & 30.00 \\
4 & 7 & 7.00 \\
5 & 4 & 4.00 \\
\bottomrule
Total & 100 & 100.00 \\
\end{tabular}
\end{table}

\begin{table}[H]
\centering
\caption{Frecventa consumului dupa sexe}
\end{table}


\begin{table}[H]
\centering
\caption{Structura după tipul de cafea preferat}
\label{tab:tip_cafea}
\begin{tabular}{@{}llllll@{}}
\toprule
Tip & Espresso & Latte & Cappuccino & Americano & Alte \\
\midrule
Frecvență &  &  &  &  &  \\
Pondere (\%) &  &  &  &  &  \\
\bottomrule
\end{tabular}
\end{table}

\subsection{Gruparea după caracteristici cantitative (intervale discrete/continue)}
Definiți intervalele folosind reguli de tip Sturges: 
$$
k = 1 + 3.322 \cdot \log_{10}(n), \quad h = \frac{x_{\max} - x_{\min}}{k}.
$$

\begin{table}[H]
\centering
\caption{Gruparea după numărul de căni pe zi (intervale continue)}
\label{tab:cani_intervale}
\begin{tabular}{@{}lllll@{}}
\toprule
Interval & [0,1) & [1,2) & [2,3) & [3,4) \\
\midrule
Frecvență &  &  &  &  \\
Pondere (\%) &  &  &  &  \\
\bottomrule
\end{tabular}
\end{table}

\subsection{Prezentare grafică a distribuțiilor și a structurii}
Includeți diagrame (histograme, bar chart, pie chart) exportate din Jupyter.

% \begin{figure}[H]
%   \centering
%   \includegraphics[width=0.7\textwidth]{figuri/hist_cani_pe_zi.png}
%   \caption{Distribuția numărului de căni pe zi}
%   \label{fig:hist_cani}
% \end{figure}

% \begin{figure}[H]
%   \centering
%   \includegraphics[width=0.65\textwidth]{figuri/bar_tip_cafea.png}
%   \caption{Structura după tipul de cafea preferat}
%   \label{fig:bar_tip_cafea}
% \end{figure}

% III. ANALIZA STATISTICĂ INFERENȚIALĂ
\section{Analiza statistică inferențială folosind indicatori de tendință centrală}

\subsection{Indicatori de medie (pentru fiecare distribuție/grupare)}
Exemple (pentru Căni/zi):
$$
\bar{x} = \frac{1}{n}\sum_{i=1}^{n} x_i, \quad 
\tilde{x}=\text{mediană}, \quad
\text{media pătratică } = \sqrt{\frac{1}{n}\sum x_i^2}.
$$
\

\begin{appendices}
\appendix
\section*{Anexa I. Setul complet de date}
\begin{longtable}{rlllrlrl}
\toprule
Vârsta & Sex & Localitate & Tip Cafea & Cani/zi & Perioada & Buget (MDL) & Loc \\
\midrule
\endfirsthead
\toprule
Vârsta & Sex & Localitate & Tip Cafea & Cani/zi & Perioada & Buget (MDL) & Loc \\
\midrule
\endhead
\midrule
\multicolumn{8}{r}{Continued on next page} \\
\midrule
\endfoot
\bottomrule
\endlastfoot
18 & F & Urban & Latte & 5 & Dimineata & 796 & Supermarket \\
42 & F & Rural & Americano & 2 & Dimineata & 497 & Oficiu \\
60 & M & Rural & Decaf & 4 & Seara & 1151 & Oficiu \\
57 & F & Urban & Cappucino & 3 & Dimineata & 276 & Supermarket \\
59 & M & Rural & Latte & 3 & Pranz & 1307 & Supermarket \\
42 & M & Urban & Cappucino & 3 & Seara & 2022 & Supermarket \\
56 & F & Urban & Decaf & 2 & Dimineata & 2407 & Oficiu \\
52 & M & Urban & Espresso & 2 & Dimineata & 2041 & Supermarket \\
20 & M & Urban & Americano & 2 & Dimineata & 1768 & Distributie \\
67 & M & Urban & Espresso & 3 & Pranz & 2025 & Distributie \\
54 & F & Rural & Americano & 1 & Pranz & 1908 & Cafenea \\
61 & M & Urban & Americano & 2 & Pranz & 1746 & Oficiu \\
61 & M & Rural & Espresso & 3 & Dimineata & 758 & Supermarket \\
68 & M & Urban & Americano & 3 & Pranz & 1362 & Oficiu \\
47 & M & Rural & Espresso & 3 & Seara & 340 & Distributie \\
55 & M & Rural & Espresso & 3 & Dimineata & 1086 & Cafenea \\
51 & M & Urban & Americano & 5 & Seara & 1678 & Oficiu \\
29 & F & Rural & Latte & 2 & Pranz & 1500 & Cafenea \\
36 & M & Urban & Americano & 2 & Pranz & 1143 & Cafenea \\
54 & M & Rural & Latte & 2 & Dimineata & 1443 & Supermarket \\
61 & M & Urban & Decaf & 2 & Dimineata & 1562 & Supermarket \\
66 & M & Rural & Espresso & 4 & Dimineata & 720 & Alte \\
34 & F & Urban & Americano & 2 & Dimineata & 265 & Supermarket \\
27 & F & Urban & Latte & 3 & Dimineata & 1412 & Supermarket \\
66 & M & Rural & Decaf & 2 & Pranz & 577 & Oficiu \\
64 & M & Urban & Decaf & 3 & Seara & 2227 & Cafenea \\
29 & M & Urban & Latte & 3 & Pranz & 2092 & Alte \\
33 & F & Urban & Americano & 2 & Dimineata & 1037 & Oficiu \\
41 & F & Urban & Cappucino & 2 & Dimineata & 962 & Distributie \\
36 & F & Urban & Decaf & 2 & Dimineata & 440 & Cafenea \\
25 & M & Urban & Latte & 1 & Pranz & 1050 & Supermarket \\
48 & F & Urban & Americano & 2 & Pranz & 1434 & Supermarket \\
38 & M & Urban & Cappucino & 1 & Seara & 1697 & Distributie \\
34 & M & Rural & Espresso & 2 & Dimineata & 1058 & Cafenea \\
40 & M & Urban & Espresso & 3 & Seara & 1668 & Supermarket \\
54 & M & Urban & Americano & 3 & Pranz & 2047 & Alte \\
33 & F & Urban & Espresso & 3 & Dimineata & 275 & Oficiu \\
23 & F & Urban & Americano & 1 & Pranz & 978 & Supermarket \\
25 & M & Rural & Espresso & 3 & Dimineata & 730 & Cafenea \\
42 & F & Urban & Latte & 2 & Dimineata & 2085 & Alte \\
35 & F & Urban & Cappucino & 2 & Dimineata & 1510 & Cafenea \\
42 & F & Rural & Decaf & 4 & Dimineata & 1390 & Distributie \\
29 & F & Rural & Espresso & 1 & Dimineata & 1425 & Distributie \\
32 & F & Urban & Americano & 3 & Dimineata & 1246 & Cafenea \\
43 & F & Rural & Cappucino & 3 & Dimineata & 345 & Distributie \\
58 & F & Rural & Espresso & 2 & Dimineata & 1114 & Cafenea \\
62 & M & Urban & Decaf & 2 & Dimineata & 2420 & Supermarket \\
43 & M & Rural & Latte & 2 & Dimineata & 1192 & Distributie \\
64 & F & Rural & Latte & 4 & Dimineata & 1662 & Acasa \\
49 & M & Urban & Decaf & 2 & Dimineata & 530 & Oficiu \\
27 & F & Urban & Espresso & 2 & Dimineata & 970 & Oficiu \\
33 & M & Urban & Americano & 2 & Dimineata & 993 & Cafenea \\
24 & M & Rural & Latte & 3 & Dimineata & 1784 & Cafenea \\
34 & M & Rural & Espresso & 1 & Dimineata & 682 & Cafenea \\
40 & M & Rural & Americano & 2 & Pranz & 1331 & Cafenea \\
43 & F & Urban & Decaf & 3 & Dimineata & 1659 & Supermarket \\
38 & F & Urban & Espresso & 4 & Dimineata & 1081 & Supermarket \\
39 & F & Urban & Cappucino & 2 & Seara & 921 & Supermarket \\
24 & F & Urban & Latte & 2 & Pranz & 1029 & Cafenea \\
31 & M & Urban & Americano & 2 & Seara & 1050 & Oficiu \\
32 & F & Rural & Americano & 3 & Pranz & 1655 & Distributie \\
24 & F & Urban & Latte & 2 & Dimineata & 1225 & Supermarket \\
26 & F & Rural & Decaf & 2 & Dimineata & 1878 & Acasa \\
68 & M & Urban & Decaf & 3 & Dimineata & 682 & Acasa \\
65 & F & Urban & Decaf & 3 & Dimineata & 1382 & Cafenea \\
25 & F & Urban & Latte & 2 & Dimineata & 2059 & Distributie \\
67 & F & Urban & Americano & 4 & Dimineata & 1562 & Oficiu \\
40 & F & Rural & Americano & 3 & Dimineata & 1746 & Oficiu \\
46 & M & Urban & Americano & 2 & Pranz & 1051 & Supermarket \\
35 & F & Rural & Americano & 1 & Pranz & 1168 & Oficiu \\
48 & F & Urban & Americano & 5 & Dimineata & 985 & Supermarket \\
47 & M & Rural & Cappucino & 1 & Pranz & 1300 & Oficiu \\
69 & F & Urban & Decaf & 2 & Dimineata & 917 & Oficiu \\
56 & F & Rural & Espresso & 1 & Pranz & 640 & Distributie \\
52 & M & Rural & Espresso & 3 & Dimineata & 1477 & Oficiu \\
35 & M & Urban & Latte & 2 & Seara & 1279 & Distributie \\
59 & M & Urban & Americano & 1 & Pranz & 1225 & Cafenea \\
56 & M & Rural & Americano & 1 & Dimineata & 1543 & Cafenea \\
34 & M & Rural & Americano & 1 & Dimineata & 462 & Distributie \\
31 & F & Urban & Americano & 3 & Dimineata & 2191 & Supermarket \\
48 & M & Rural & Americano & 3 & Dimineata & 1881 & Cafenea \\
41 & M & Rural & Latte & 2 & Dimineata & 1777 & Cafenea \\
52 & F & Rural & Americano & 2 & Dimineata & 1597 & Cafenea \\
61 & F & Rural & Decaf & 3 & Dimineata & 836 & Alte \\
62 & M & Rural & Cappucino & 1 & Pranz & 1390 & Oficiu \\
51 & M & Urban & Latte & 2 & Seara & 2321 & Acasa \\
20 & M & Urban & Espresso & 5 & Dimineata & 739 & Oficiu \\
54 & F & Urban & Espresso & 3 & Dimineata & 1494 & Oficiu \\
60 & M & Rural & Espresso & 2 & Dimineata & 1503 & Supermarket \\
57 & F & Rural & Americano & 2 & Pranz & 899 & Oficiu \\
43 & F & Urban & Latte & 2 & Dimineata & 574 & Cafenea \\
40 & M & Rural & Americano & 2 & Seara & 1754 & Acasa \\
61 & F & Rural & Decaf & 1 & Pranz & 1803 & Distributie \\
56 & F & Urban & Decaf & 1 & Dimineata & 1692 & Supermarket \\
32 & M & Urban & Latte & 1 & Dimineata & 1474 & Alte \\
66 & F & Urban & Espresso & 4 & Pranz & 948 & Distributie \\
21 & F & Urban & Latte & 3 & Pranz & 610 & Distributie \\
46 & F & Urban & Espresso & 1 & Dimineata & 1732 & Cafenea \\
39 & F & Urban & Cappucino & 1 & Dimineata & 951 & Oficiu \\
42 & F & Urban & Espresso & 3 & Pranz & 1367 & Oficiu \\
\end{longtable}

\end{appendices}

\end{document}